\section{Society Simulation}~\label{sec:soc}


\begin{table*}[!htp]
\setlength{\tabcolsep}{2pt}
\centering
\resizebox*{!}{\dimexpr\textheight-2\baselineskip\relax}{
\begin{tabular}{c|c|c|c|cccc}
\toprule
\multirow{2}{*}{\textbf{Scenario}}                                                            & \multirow{2}{*}{\textbf{Field}}                                                                                     & \multirow{2}{*}{\textbf{Paper}}                            & \multirow{2}{*}{\textbf{\# Agents}} & \multicolumn{4}{c}{\textbf{Construction Element}}                                                                                                     \\ \cline{5-8} 
                                                                                              &                                                                                                                     &                                                            &                                     & \multicolumn{1}{c|}{\textbf{Composition}} & \multicolumn{1}{c|}{\textbf{Network}} & \multicolumn{1}{c|}{\textbf{Social Influence}} & \textbf{Outcome} \\ \midrule
\multirow{17}{*}{\begin{tabular}[c]{@{}c@{}}General\\ Economic\end{tabular}}                  & \multirow{10}{*}{\begin{tabular}[c]{@{}c@{}}Game Theory\\ and\\ Strategic Interactions\end{tabular}}                & Agent-trust~\cite{xie2024can}                              & $(0, 10]$                           & \multicolumn{1}{c|}{\checkmark}           & \multicolumn{1}{c|}{}                 & \multicolumn{1}{c|}{\checkmark}                & \checkmark       \\ \cline{3-8} 
                                                                                              &                                                                                                                     & LELMA~\cite{mensfelt2024logic}                             & $(0, 10]$                           & \multicolumn{1}{c|}{}                     & \multicolumn{1}{c|}{}                 & \multicolumn{1}{c|}{\checkmark}                & \checkmark       \\ \cline{3-8} 
                                                                                              &                                                                                                                     & Economics   Arena~\cite{guo2024economics}                  & $(0, 10]$                           & \multicolumn{1}{c|}{}                     & \multicolumn{1}{c|}{}                 & \multicolumn{1}{c|}{\checkmark}                & \checkmark       \\ \cline{3-8} 
                                                                                              &                                                                                                                     & Fontana et   al.~\cite{fontana2024nicer}                   & $(0, 10]$                           & \multicolumn{1}{c|}{}                     & \multicolumn{1}{c|}{}                 & \multicolumn{1}{c|}{\checkmark}                & \checkmark       \\ \cline{3-8} 
                                                                                              &                                                                                                                     & SABM~\cite{han2023guinea}                                  & $(0, 10]$                           & \multicolumn{1}{c|}{\checkmark}           & \multicolumn{1}{c|}{\checkmark}       & \multicolumn{1}{c|}{\checkmark}                & \checkmark       \\ \cline{3-8} 
                                                                                              &                                                                                                                     & Noh and   Chang.~\cite{noh2024llms}                        & $(0, 10]$                           & \multicolumn{1}{c|}{\checkmark}           & \multicolumn{1}{c|}{}                 & \multicolumn{1}{c|}{\checkmark}                & \checkmark       \\ \cline{3-8} 
                                                                                              &                                                                                                                     & Mozikov et   al.~\cite{mozikov2024good}                    & $(0, 10]$                           & \multicolumn{1}{c|}{}                     & \multicolumn{1}{c|}{}                 & \multicolumn{1}{c|}{\checkmark}                & \checkmark       \\ \cline{3-8} 
                                                                                              &                                                                                                                     & Wu et al.~\cite{wu2024shall}                               & $(10, 100]$                         & \multicolumn{1}{c|}{}                     & \multicolumn{1}{c|}{\checkmark}       & \multicolumn{1}{c|}{\checkmark}                & \checkmark       \\ \cline{3-8} 
                                                                                              &                                                                                                                     & CompeteAI~\cite{zhaocompeteai}                             & $(10, 100]$                         & \multicolumn{1}{c|}{\checkmark}           & \multicolumn{1}{c|}{\checkmark}       & \multicolumn{1}{c|}{\checkmark}                & \checkmark       \\ \cline{3-8} 
                                                                                              &                                                                                                                     & WarAgent~\cite{hua2023war}                                 & $(10, 100]$                         & \multicolumn{1}{c|}{\checkmark}           & \multicolumn{1}{c|}{\checkmark}       & \multicolumn{1}{c|}{\checkmark}                & \checkmark       \\ \cline{2-8} 
                                                                                              & \multirow{7}{*}{\begin{tabular}[c]{@{}c@{}}Economic \\ Contexts\end{tabular}}                                       & Horton~\cite{horton2023large}                              & $(10, 100]$                         & \multicolumn{1}{c|}{\checkmark}           & \multicolumn{1}{c|}{}                 & \multicolumn{1}{c|}{}                          & \checkmark       \\ \cline{3-8} 
                                                                                              &                                                                                                                     & EconAgent~\cite{li2024econagent}                           & $(10, 100]$                         & \multicolumn{1}{c|}{\checkmark}           & \multicolumn{1}{c|}{}                 & \multicolumn{1}{c|}{}                          & \checkmark       \\ \cline{3-8} 
                                                                                              &                                                                                                                     & SRAP-Agent~\cite{ji2024srap}                               & $(10, 100]$                         & \multicolumn{1}{c|}{\checkmark}           & \multicolumn{1}{c|}{\checkmark}       & \multicolumn{1}{c|}{\checkmark}                & \checkmark       \\ \cline{3-8} 
                                                                                              &                                                                                                                     & Ghaffarzadegan et   al.\cite{ghaffarzadegan2023generative} & $(10, 100]$                         & \multicolumn{1}{c|}{\checkmark}           & \multicolumn{1}{c|}{}                 & \multicolumn{1}{c|}{\checkmark}                & \checkmark       \\ \cline{3-8} 
                                                                                              &                                                                                                                     & EC~\cite{de2023emergent}                                   & $(10, 100]$                         & \multicolumn{1}{c|}{\checkmark}           & \multicolumn{1}{c|}{\checkmark}       & \multicolumn{1}{c|}{\checkmark}                & \checkmark       \\ \cline{3-8} 
                                                                                              &                                                                                                                     & Williams et   al.~\cite{williams2023epidemic}              & $(100, \infty)$                     & \multicolumn{1}{c|}{\checkmark}           & \multicolumn{1}{c|}{\checkmark}       & \multicolumn{1}{c|}{\checkmark}                & \checkmark       \\ \cline{3-8} 
                                                                                              &                                                                                                                     & AgentTorch~\cite{chopra2024limits}                         & $(100, \infty)$                     & \multicolumn{1}{c|}{\checkmark}           & \multicolumn{1}{c|}{\checkmark}       & \multicolumn{1}{c|}{}                          & \checkmark       \\ \hline
\multirow{24}{*}{\begin{tabular}[c]{@{}c@{}}Sociology\\ and\\ Political Science\end{tabular}} & \multirow{7}{*}{\begin{tabular}[c]{@{}c@{}}Public Opinion\\ Survey\end{tabular}}                                    & Argyle et al.~\cite{Argyle_2023}                           & $(100, \infty)$                     & \multicolumn{1}{c|}{\checkmark}           & \multicolumn{1}{c|}{}                 & \multicolumn{1}{c|}{}                          & \checkmark       \\ \cline{3-8} 
                                                                                              &                                                                                                                     & Lee et al.~\cite{lee2023can}                               & $(100, \infty)$                     & \multicolumn{1}{c|}{\checkmark}           & \multicolumn{1}{c|}{}                 & \multicolumn{1}{c|}{}                          & \checkmark       \\ \cline{3-8} 
                                                                                              &                                                                                                                     & Chaudhary and   Chaudhary~\cite{chaudhary2024large}        & $(100, \infty)$                     & \multicolumn{1}{c|}{\checkmark}           & \multicolumn{1}{c|}{}                 & \multicolumn{1}{c|}{}                          & \checkmark       \\ \cline{3-8} 
                                                                                              &                                                                                                                     & ElectionSim~\cite{zhang2024electionsim}                    & $(100, \infty)$                     & \multicolumn{1}{c|}{\checkmark}           & \multicolumn{1}{c|}{}                 & \multicolumn{1}{c|}{}                          & \checkmark       \\ \cline{3-8} 
                                                                                              &                                                                                                                     & GABSS~\cite{xiao2023simulating}                            & $(100, \infty)$                     & \multicolumn{1}{c|}{\checkmark}           & \multicolumn{1}{c|}{\checkmark}       & \multicolumn{1}{c|}{\checkmark}                & \checkmark       \\ \cline{3-8} 
                                                                                              &                                                                                                                     & Park et   al.~\cite{park2024generative}                    & $(100, \infty)$                     & \multicolumn{1}{c|}{\checkmark}           & \multicolumn{1}{c|}{}                 & \multicolumn{1}{c|}{}                          & \checkmark       \\ \cline{3-8} 
                                                                                              &                                                                                                                     & Sun et   al.~\cite{sun2024randomsiliconsamplingsimulating} & $(100, \infty)$                     & \multicolumn{1}{c|}{\checkmark}           & \multicolumn{1}{c|}{}                 & \multicolumn{1}{c|}{}                          & \checkmark       \\ \cline{2-8} 
                                                                                              & \multirow{17}{*}{\begin{tabular}[c]{@{}c@{}}Individual \\ and\\ Organizational\\ Behavior Observation\end{tabular}} & Aher et al.~\cite{aher2023using}                           & $(0, 10]$                           & \multicolumn{1}{c|}{}                     & \multicolumn{1}{c|}{}                 & \multicolumn{1}{c|}{\checkmark}                & \checkmark       \\ \cline{3-8} 
                                                                                              &                                                                                                                     & Zhang et   al.~\cite{zhang2023exploring}                   & $(0, 10]$                           & \multicolumn{1}{c|}{\checkmark}           & \multicolumn{1}{c|}{}                 & \multicolumn{1}{c|}{}                          & \checkmark       \\ \cline{3-8} 
                                                                                              &                                                                                                                     & Lyfe   Agents~\cite{kaiya2023lyfe}                         & $(0, 10]$                           & \multicolumn{1}{c|}{\checkmark}           & \multicolumn{1}{c|}{\checkmark}       & \multicolumn{1}{c|}{\checkmark}                & \checkmark       \\ \cline{3-8} 
                                                                                              &                                                                                                                     & CRSEC~\cite{ren2024emergence}                              & $(0, 10]$                           & \multicolumn{1}{c|}{\checkmark}           & \multicolumn{1}{c|}{\checkmark}       & \multicolumn{1}{c|}{\checkmark}                & \checkmark       \\ \cline{3-8} 
                                                                                              &                                                                                                                     & Chuang et   al.\cite{chuang2023simulating}                 & $(0, 10]$                           & \multicolumn{1}{c|}{\checkmark}           & \multicolumn{1}{c|}{\checkmark}       & \multicolumn{1}{c|}{\checkmark}                & \checkmark       \\ \cline{3-8} 
                                                                                              &                                                                                                                     & ChoiceMates~\cite{park2023choicemates}                     & $(0, 10]$                           & \multicolumn{1}{c|}{\checkmark}           & \multicolumn{1}{c|}{\checkmark}       & \multicolumn{1}{c|}{\checkmark}                & \checkmark       \\ \cline{3-8} 
                                                                                              &                                                                                                                     & Jarrett et   al.\cite{jarrett2023language}                 & $(0, 10]$                           & \multicolumn{1}{c|}{\checkmark}           & \multicolumn{1}{c|}{}                 & \multicolumn{1}{c|}{}                          & \checkmark       \\ \cline{3-8} 
                                                                                              &                                                                                                                     & AgentReview~\cite{jin2024agentreview}                      & $(0, 10]$                           & \multicolumn{1}{c|}{\checkmark}           & \multicolumn{1}{c|}{\checkmark}       & \multicolumn{1}{c|}{}                          & \checkmark       \\ \cline{3-8} 
                                                                                              &                                                                                                                     & Generative   Agents~\cite{park2023generative}              & $(10, 100]$                         & \multicolumn{1}{c|}{\checkmark}           & \multicolumn{1}{c|}{\checkmark}       & \multicolumn{1}{c|}{\checkmark}                & \checkmark       \\ \cline{3-8} 
                                                                                              &                                                                                                                     & AGA~\cite{yu2024affordable}                                & $(10, 100]$                         & \multicolumn{1}{c|}{\checkmark}           & \multicolumn{1}{c|}{\checkmark}       & \multicolumn{1}{c|}{\checkmark}                & \checkmark       \\ \cline{3-8} 
                                                                                              &                                                                                                                     & MineLand~\cite{yu2024mineland}                             & $(10, 100]$                         & \multicolumn{1}{c|}{\checkmark}           & \multicolumn{1}{c|}{\checkmark}       & \multicolumn{1}{c|}{\checkmark}                & \checkmark       \\ \cline{3-8} 
                                                                                              &                                                                                                                     & Chuang et   al.~\cite{chuang2024wisdom}                    & $(10, 100]$                         & \multicolumn{1}{c|}{\checkmark}           & \multicolumn{1}{c|}{\checkmark}       & \multicolumn{1}{c|}{\checkmark}                & \checkmark       \\ \cline{3-8} 
                                                                                              &                                                                                                                     & CareerAgent~\cite{zhu2024generative}                       & $(10, 100]$                         & \multicolumn{1}{c|}{\checkmark}           & \multicolumn{1}{c|}{\checkmark}       & \multicolumn{1}{c|}{\checkmark}                & \checkmark       \\ \cline{3-8} 
                                                                                              &                                                                                                                     & Suzuki and   Arita~\cite{suzuki2024evolutionary}           & $(10, 100]$                         & \multicolumn{1}{c|}{\checkmark}           & \multicolumn{1}{c|}{}                 & \multicolumn{1}{c|}{\checkmark}                & \checkmark       \\ \cline{3-8} 
                                                                                              &                                                                                                                     & Chuang et   al.\cite{chuang2024beyond}                     & $(100, \infty)$                     & \multicolumn{1}{c|}{\checkmark}           & \multicolumn{1}{c|}{}                 & \multicolumn{1}{c|}{}                          & \checkmark       \\ \cline{3-8} 
                                                                                              &                                                                                                                     & Li et   al.~\cite{li2023quantifying}                       & $(100, \infty)$                     & \multicolumn{1}{c|}{}                     & \multicolumn{1}{c|}{\checkmark}       & \multicolumn{1}{c|}{\checkmark}                & \checkmark       \\ \cline{3-8} 
                                                                                              &                                                                                                                     & MATRIX~\cite{tang2024synthesizing}                         & $(100, \infty)$                     & \multicolumn{1}{c|}{\checkmark}           & \multicolumn{1}{c|}{\checkmark}       & \multicolumn{1}{c|}{}                          &                  \\ \hline
\multirow{19}{*}{\begin{tabular}[c]{@{}c@{}}Online\\ Platform\end{tabular}}                   & \multirow{14}{*}{\begin{tabular}[c]{@{}c@{}}Social\\ Platforms\end{tabular}}                                        & Cai et al.\cite{cai2024language}                           & $(0, 10]$                           & \multicolumn{1}{c|}{}                     & \multicolumn{1}{c|}{}                 & \multicolumn{1}{c|}{\checkmark}                & \checkmark       \\ \cline{3-8} 
                                                                                              &                                                                                                                     & FPS~\cite{liu2024skepticism}                               & $(10, 100]$                         & \multicolumn{1}{c|}{\checkmark}           & \multicolumn{1}{c|}{\checkmark}       & \multicolumn{1}{c|}{\checkmark}                & \checkmark       \\ \cline{3-8} 
                                                                                              &                                                                                                                     & FUSE~\cite{liu2024tiny}                                    & $(10, 100]$                         & \multicolumn{1}{c|}{\checkmark}           & \multicolumn{1}{c|}{\checkmark}       & \multicolumn{1}{c|}{\checkmark}                & \checkmark       \\ \cline{3-8} 
                                                                                              &                                                                                                                     & Wang et   al.\cite{wang2024decoding}                       & $(10, 100]$                         & \multicolumn{1}{c|}{\checkmark}           & \multicolumn{1}{c|}{\checkmark}       & \multicolumn{1}{c|}{\checkmark}                & \checkmark       \\ \cline{3-8} 
                                                                                              &                                                                                                                     & Concordia~\cite{touzel2024simulation}                      & $(10, 100]$                         & \multicolumn{1}{c|}{\checkmark}           & \multicolumn{1}{c|}{\checkmark}       & \multicolumn{1}{c|}{\checkmark}                & \checkmark       \\ \cline{3-8} 
                                                                                              &                                                                                                                     & Social   Simulacra~\cite{park2022social}                   & $(100, \infty)$                     & \multicolumn{1}{c|}{\checkmark}           & \multicolumn{1}{c|}{\checkmark}       & \multicolumn{1}{c|}{\checkmark}                & \checkmark       \\ \cline{3-8} 
                                                                                              &                                                                                                                     & $S^3$~\cite{gao2023s3}                                     & $(100, \infty)$                     & \multicolumn{1}{c|}{\checkmark}           & \multicolumn{1}{c|}{\checkmark}       & \multicolumn{1}{c|}{\checkmark}                & \checkmark       \\ \cline{3-8} 
                                                                                              &                                                                                                                     & Törnberg et   al.~\cite{tornberg2023simulating}            & $(100, \infty)$                     & \multicolumn{1}{c|}{\checkmark}           & \multicolumn{1}{c|}{\checkmark}       & \multicolumn{1}{c|}{\checkmark}                & \checkmark       \\ \cline{3-8} 
                                                                                              &                                                                                                                     & Y   Social~\cite{rossetti2024social}                       & $(100, \infty)$                     & \multicolumn{1}{c|}{\checkmark}           & \multicolumn{1}{c|}{\checkmark}       & \multicolumn{1}{c|}{\checkmark}                & \checkmark       \\ \cline{3-8} 
                                                                                              &                                                                                                                     & TIS~\cite{zhang2024large}                                  & $(100, \infty)$                     & \multicolumn{1}{c|}{\checkmark}           & \multicolumn{1}{c|}{\checkmark}       & \multicolumn{1}{c|}{\checkmark}                & \checkmark       \\ \cline{3-8} 
                                                                                              &                                                                                                                     & HiSim~\cite{mou2024unveiling}                              & $(100, \infty)$                     & \multicolumn{1}{c|}{\checkmark}           & \multicolumn{1}{c|}{\checkmark}       & \multicolumn{1}{c|}{\checkmark}                & \checkmark       \\ \cline{3-8} 
                                                                                              &                                                                                                                     & OASIS~\cite{yang2024oasis}                                 & $(100, \infty)$                     & \multicolumn{1}{c|}{\checkmark}           & \multicolumn{1}{c|}{\checkmark}       & \multicolumn{1}{c|}{\checkmark}                & \checkmark       \\ \cline{3-8} 
                                                                                              &                                                                                                                     & MindEcho~\cite{xu2024mindechoroleplayinglanguageagents}    & $(100, \infty)$                     & \multicolumn{1}{c|}{\checkmark}           & \multicolumn{1}{c|}{}                 & \multicolumn{1}{c|}{\checkmark}                &                  \\ \cline{3-8} 
                                                                                              &                                                                                                                     & BASES~\cite{ren2024bases}                                  & $(100, \infty)$                     & \multicolumn{1}{c|}{\checkmark}           & \multicolumn{1}{c|}{}                 & \multicolumn{1}{c|}{}                          &                  \\ \cline{2-8} 
                                                                                              & \multirow{5}{*}{\begin{tabular}[c]{@{}c@{}}Recommendation\\ Environments\end{tabular}}                              & InteRecAgent~\cite{huang2023recommender}                   & $(0, 10]$                           & \multicolumn{1}{c|}{\checkmark}           & \multicolumn{1}{c|}{}                 & \multicolumn{1}{c|}{}                          &                  \\ \cline{3-8} 
                                                                                              &                                                                                                                     & Rec4Agentverse~\cite{zhang2024prospect}                    & $(0, 10]$                           & \multicolumn{1}{c|}{\checkmark}           & \multicolumn{1}{c|}{\checkmark}       & \multicolumn{1}{c|}{}                          &                  \\ \cline{3-8} 
                                                                                              &                                                                                                                     & RecAgent~\cite{wang2023recagent}                           & $(10, 100]$                         & \multicolumn{1}{c|}{\checkmark}           & \multicolumn{1}{c|}{\checkmark}       & \multicolumn{1}{c|}{\checkmark}                & \checkmark       \\ \cline{3-8} 
                                                                                              &                                                                                                                     & Agent4Rec~\cite{Zhang2024on}                               & $(100, \infty)$                     & \multicolumn{1}{c|}{\checkmark}           & \multicolumn{1}{c|}{\checkmark}       & \multicolumn{1}{c|}{\checkmark}                & \checkmark       \\ \cline{3-8} 
                                                                                              &                                                                                                                     & AgentCF~\cite{zhang2024agentcf}                            & $(100, \infty)$                     & \multicolumn{1}{c|}{\checkmark}           & \multicolumn{1}{c|}{\checkmark}       & \multicolumn{1}{c|}{\checkmark}                & \checkmark       \\ \bottomrule
\end{tabular}
}
\caption{A list of representative works of society simulation.}
\label{tab:soc}
\end{table*}


While scenarios discuss multi-agent interactions in relatively focused and small-scale contexts and provide solutions within specific domains, society is more complex than a simple scenario. Its complexity lies in many aspects, such as the diversity of its components, the variety of structures, and nonlinear effects~\cite{squazzoni2014social}. Considering this, a series of studies focus on society simulation. In terms of research topic, society simulation generally hopes to investigate societal and macro-level results. In terms of research purpose, society simulation does not aim to solve a task or problem, instead, it focuses on revealing and explaining emergent behaviors and the outcomes of interactions among numerous agents. Society simulations have been a vital tool for theoretical validation and predicting social dynamics.

In this section, we summarize the components of social construction to capture the key features reflected in society simulations in~\S\ref{subsec:soc_element}. Then, we present the different categories of scenarios in society simulation in~\S\ref{subsec:soc_scenario}. After that, we introduce the evaluation of society simulation in~\S\ref{subsec:soc_eval}. The overall framework is illustrated in Figure~\ref{fig:fm_soc} and representative works are summarized in Table~\ref{tab:soc}.


\subsection{Social Construction Elements}
\label{subsec:soc_element}
Considering the complexity of society, a major challenge in society simulation is bridging the gap between individual and societal scales. Some core elements serve as the foundation for modeling social systems. We outline four key dimensions that underpin societal structures and dynamics: \textbf{composition}, \textbf{network},\textbf{ social influence}, and \textbf{outcomes}.




\subsubsection{Composition}
Society is composed of massive and diverse individuals. This diversity, also referred to as heterogeneity~\cite{squazzoni2014social} in social science, encompasses a wide range of beliefs, preferences, behaviors, normative values, and positions within social structures. Modeling this diversity is essential for capturing the varied behavioral patterns and complex social dynamics that emerge from individual differences within a social system.

\paragraph{Individual Composition} 
To model a diverse society, the composition of individuals in society needs to be determined. There are three main approaches to determining the composition of individuals in a system simulating a microcosm of society. Some works rely on \textit{{\ul virtual individual synthesis}}, often not focused on alignment with the real world, aiming to ensure that the system includes users with a variety of attributes, typically by generating virtual individuals with the help of LLMs or humans~\cite{chuang2024wisdom,binz2023turning}. Other works utilize \textit{{\ul existing datasets}}, such as MovieLens-1M~\cite{wang2023recagent,Zhang2024on}, to define user composition within a simulated recommendation platform. Agents are initialized on the basis of the user information within these datasets, reflecting the distribution of users in that context. Recently, an increasing number of studies have focused on \textit{{\ul real-world distribution replication}}, such as the composition of users on social platforms~\cite{yang2024oasis} or the distribution of voters in surveys~\cite{zhang2024electionsim}. For small-scale individual sets, individual data are typically collected manually~\cite{park2023choicemates,park2024generative}. In cases where large-scale populations are required or obtaining real data is difficult, individuals may be sampled based on real-world macro distributions or generated by LLMs to match desired attribute distribution~\cite{Argyle_2023,lee2023can,zhang2024electionsim}.

\paragraph{Trade-off between Simulation Precision and Scale}
When simulating individuals in society simulations, many studies adopt detailed role modeling to enhance the authenticity of agent behavior. Beyond common demographic attributes, this may include factors such as an individual’s past statements and interaction history~\cite{park2023generative,wang2023recagent,Zhang2024on,zhaocompeteai,fontana2024nicer}. However, as the number of individuals increases, such fine-grained modeling becomes expensive. Consequently, a trade-off often arises between the precision of individual modeling and the scale of the simulation. In large-scale simulations, to reduce computational costs, the details of each agent are typically simplified, by retaining only the most essential and common characteristics~\cite{williams2023epidemic,chopra2024limits} or compressing auxiliary dialogue information into shared memory~\cite{yu2024affordable}.

\paragraph{Special Modeling on Outliers} 
As previously mentioned, the composition of individuals in society is diverse. However, not all individuals play an equally significant role. Some individuals, whose attributes or behaviors significantly deviate from the majority, are referred to as outliers~\cite{squazzoni2014social}. Compared to average individuals, outliers often introduce variability and unpredictability to society. Examples include celebrities and opinion leaders~\cite{zhang2024large,xu2024mindechoroleplayinglanguageagents}, who frequently hold prominent positions within social structures and amplify their influence. In situations with limited resources, some studies~\cite{mou2024unveiling} prioritize detailed modeling of these core content producers, while simplifying the modeling for the majority. Meanwhile, intervention policies based on simulation results often focus on these key nodes in networks~\cite{li2024large}, aiming to influence the overall system's behavior by blocking or interfering with them.


\subsubsection{Network}
Social interactions are often conducted through social networks, which can be described using graph structures where nodes represent individuals and edges represent their relations. The network determines the direction of information and influence dissemination. In social science, it has been observed that homophily of individuals can increase the likelihood of communication. Highly similar individuals are more likely to establish connections compared to those with greater differences~\cite{brown1987social,kossinets2009origins}. This principle also informs the construction of networks in society simulations. The methods for constructing social networks vary across different scenarios. Here, we divide them into offline networks and online networks.

\paragraph{Offline Network}
An offline network represents connections formed through in-person interactions, such as face-to-face communication or the spread of opinions and diseases in physical settings. On the one hand, some studies aim to simulate interactions in virtual worlds, thus determining the connections between agents in a \textit{{\ul random or predefined}} manner~\cite{park2023generative,yu2024affordable,ren2024emergence}. On the other hand, when some studies aim to simulate the spread of a disease or event information in the real world, considering the difficulty of obtaining real data, they often estimate the social relations using \textit{{\ul external algorithms}} or agents themselves~\cite{xiao2023simulating,williams2023epidemic}. However, in studies with a large scale of agents, the network relationships between individuals are sometimes ignored, and individuals are treated as independent~\cite{zhang2024electionsim}. Alternatively, some studies provide rough information, such as community statistics, in place of specific details about the agents' neighbors~\cite{chopra2024limits}.


\paragraph{Online Network}
An online network is a digital structure where individuals or entities interact through platforms, such as online social platforms and recommendation platforms, forming connections based on activities, relationships, or shared interests. At the beginning, some studies \textit{{\ul randomly}} initialize the social relations for users existing datasets~\cite{wang2023recagent} or synthesized users~\cite{liu2024skepticism}, while other efforts have focused on crawling \textit{{\ul authentic}} social relationships from social media platforms like Weibo~\cite{gao2023s3} and Twitter~\cite{mou2024unveiling}. However, as the scale of individuals increase, it may be challenging to obtain all of their authentic relationships. Therefore, some studies construct networks using a small portion of real relationship data combined with a large amount of \textit{{\ul synthetic}} relationship data~\cite{yang2024oasis}, or connect similar users based on the assumption of homophily~\cite{tang2024synthesizing}.


\subsubsection{Social Influence}
Social influence refers to the influence agents have on others and the influence they receive from others during interactions. This is also known as embeddedness in social sciences~\cite{squazzoni2014social}, which suggests that individuals behavior and decisions are influenced by their environment. When conducting society simulations, it is necessary to consider the modeling of such social influence.

\paragraph{Influence Received by the Influencee}
% fps; hisim; chuang
The same information may produce different effects when received by individuals with different traits. Currently, most studies have modeled how the influence received by the recipient varies based on their profile~\cite{gao2023s3,liu2024skepticism,yang2024oasis}. This can be easily achieved by integrating the individual's profile, memory and the information received from others into the same context. Building this, a few works further induce additional mechanisms such as cognitive bias~\cite{chuang2023simulating} and reflection on norms~\cite{ren2024emergence} to enhance agents' understanding and perception of the received messages.

\paragraph{Influence Exerted by the Influencer}
The same message conveyed by different individuals can result in varying social impacts. The Pareto distribution and the Matthew Effect~\cite{wang2023recagent,mou2024unveiling} indicate that information, influence, or attention tends to concentrate on a small group of individuals who are already dominant in the community. Therefore, when simulating social interactions, the identity, status, and reputation of the information sender are also crucial. Some studies start with real-world data to conduct detailed modeling of opinion leaders~\cite{xu2024mindechoroleplayinglanguageagents,zhang2024large}. Other studies, instead of focusing on the role of the influencer, model the influence exerted by the influencer by incorporating the relation information such as social impression memory~\cite{yu2024affordable} and share party affiliation~\cite{chuang2024wisdom}. In addition to the influence exerted by individuals, research has found that as group size increases, the impact of a single influencer may diminish. However, the influence of the group on individuals often drives them to align their behavior with the group, leading to the emergence of the herd effect~\cite{yang2024oasis}.


\subsubsection{Outcomes}
Social emergence suggests that the collective behaviors or phenomena arise from individual interactions are not a linear sum of individual actions but rather complex patterns emerge from the interactions~\cite{schelling1971dynamic,squazzoni2014social}. 
% Social emergence suggests that the collective behaviors or phenomena arising from society simulation. This is a complex pattern that is difficult to capture on a small-scale simulation~\cite{schelling1971dynamic,squazzoni2014social}.
These interaction outcomes may be measurable macro results, such as voting results and public opinion levels, or they may also be qualitative social phenomena and norms. Next, we will discuss these two types of outcomes separately.

\paragraph{Macro Statistical Results}
Macro statistical results are typically the focus of existing studies, as they are closely related to predefined research objectives such as market research, election predictions, and public opinion forecasting. These studies often aim to calculate the sum or average of the choices or opinions of all agents in the system. To get a static opinion distribution, some studies overlook the social interactions and instead directly sum up individual choices to obtain macro outcomes~\cite{sun2024randomsiliconsamplingsimulating,zhang2024electionsim}, simplifying the complexity of social dynamics. Another line of research focuses on the change of indicators by modeling multiple rounds of interactions among the agents over a period of time and then statistically analyzing the results~\cite{gao2023s3,tornberg2023simulating,han2023guinea,li2024econagent,wu2024shall}.


\paragraph{Formation of Social Phenomena and Social Norms}
In addition to the quantifiable macro results, some social phenomena and social norms are also important outcomes of social interactions. On the one hand, some studies have identified the bubble effect in recommendation systems~\cite{Zhang2024on}, echo chambers in social media~\cite{mou2024unveiling,yang2024oasis,wang2024decoding},  Matthew effect in competitive agent interactions~\cite{zhaocompeteai}, and spontaneous cooperation of competing agents~\cite{wu2024shall} by calculating additional metrics or observing the trends of primary indicators. On the other hand, some studies examine social norms as an important byproduct of social interactions. This includes simulating and testing whether community rules can shape desired social norms~\cite{park2022social}, constructing normative architecture to observe the emergence of social norms~\cite{ren2024emergence}, studying how social media language evolves in the presence of regulatory constraints~\cite{cai2024language}, and observing changes in social norms in real-world scenarios such as autonomous driving~\cite{wang2024can}.









\subsection{Scenario}
\label{subsec:soc_scenario}
Society simulation has been widely applied to various scenarios related to human society. These scenarios cover different aspects of daily human life, and existing studies can be categorized into three primary areas: \textbf{general economics}, \textbf{sociology and political science}, as well as \textbf{online platforms}. 

\subsubsection{General Economics}
% In social simulation, concepts from economics and game theory play significant roles in modeling societal systems, particularly when examining macroeconomic behavior or individual decision-making in resource allocation and competition. These studies mainly focus on how agents make decisions based on economic incentives, market rules, and resource constraints, while also exploring how interactions between groups shape overall economic trends.
Simulations in general economics analyze decision-making and behaviors related to resource allocation and competition. These studies primarily investigate how agents make decisions influenced by economic incentives, market rules and resource constraints, while also examining how interactions among groups shape broader economic trends.

\paragraph{Game Theory and Strategic Interactions}
Some research mainly focuses on game theory and strategic interaction. These scenarios typically involve small groups of agents, with a primary focus on the complex interactions between agents. Some works use classic game theory games, such as the Prisoner's Dilemma, to explore agent behavior in game-theoretic scenarios, including trust behavior~\cite{xie2024can}, logic reasoning and decision-making~\cite{mensfelt2024logic}, rationality and strategic reasoning ability~\cite{guo2024economics}, cooperation tendencies~\cite{fontana2024nicer} and how emotional states can disrupt rational decision-making~\cite{mozikov2024good}. Other studies focus on real-world scenarios other than the games, such as spontaneous cooperation in competitive environments~\cite{wu2024shall}, complex market behaviors in firm competition~\cite{han2023guinea}, and competition between restaurant and customer agents~\cite{zhaocompeteai}. Overall, the former kind of scenarios simplifies the environment, making it easier to conduct controlled research on agent behavior, while the latter provides more insights for real-world applications.

\paragraph{Economic Contexts}
In addition to close studies on game theory and strategic interactions, some studies focus on the use of agents and their interactions within economic environments. Horton~\cite{horton2023large} examines economic agents driven by LLMs in various experiments to replicate human behavior in economic scenarios. EconAgent~\cite{li2024econagent} introduces agents for macroeconomic simulation, emphasizing the influence of macroeconomic trends. SRAP-Agent~\cite{ji2024srap} proposes a framework for simulating and optimizing scarce resource allocation in economics, specifically in public housing allocation scenarios. Besides, some studies involve broader macroeconomic domains, using agents to simulate and predict the spread of diseases and the change in unemployment rates~\cite{williams2023epidemic,chopra2024limits}. 





\subsubsection{Sociology and Political Science}
Society simulation has been widely used in sociological and political science research.  These studies range from small-scale laboratory experiments that validate theories and hypotheses to large-scale social surveys aimed at understanding public choices. The goal is to leverage agents as substitutes for humans in studying human behavior within sociological and political contexts.


\paragraph{Public Opinion Survey} A mainstream application of society simulation is public opinion survey, which aims to predict the perspectives of specific groups toward a given subject through simulation and aggregate their opinions to support advanced needs such as election forecasting and public administration. Argyle et al.~\cite{Argyle_2023} first propose that LLMs could serve as silicon samples of humans, through several large-scale surveys conducted in the United States. Building on this, some studies have expanded their focus to scenarios of opinion surveys~\cite{lee2023can,chaudhary2024large,chuang2024beyond}, such as election polls~\cite{zhang2024electionsim} and response to public administration crisis~\cite{xiao2023simulating}, delving deeper into issues like population complexity and algorithmic bias. Recently, agents have demonstrated the potential to replicate participants' responses in individual interviews~\cite{park2024generative}. These studies lay the foundation for new tools to investigate individual and collective behavior.


\paragraph{Individual and Organizational Behavior Observation}
Other studies focus on observing individual or organizational behavior in common or specific settings. Some works do not specify a particular scenario but instead observe agents' social interactions and potential phenomena in daily life within a sandbox environment~\cite{park2023generative,kaiya2023lyfe,ren2024emergence,yu2024mineland}. 
Other studies aim to validate theories or hypotheses in specific scenarios, such as the wisdom of partisan crowds~\cite{chuang2024wisdom}, information management~\cite{park2023choicemates}, organizational behavior management~\cite{zhu2024generative}, and the evolution of personality traits~\cite{suzuki2024evolutionary}.



\subsubsection{Online Platform}
Online Platforms are a vital component of society simulation, offering a practical means to study complex social phenomena in digital environments. These platforms, ranging from social media to online communities, allow agents to simulate real-world interactions and study dynamics such as opinion formation, information spread, and collective behaviors. 
\paragraph{Social Platforms}
Online social platforms have long served as an important testing ground for studying the propagation of information and the evolution of opinions. These studies typically recreate environments similar to popular social platforms, such as Twitter, Reddit, and Weibo, with action spaces that include behaviors like sharing, commenting, and liking. By simulating these scenarios, researchers can model the spread of information and track changes in user attitudes following events, covering a wide range of topics such as general news, rumors, and the role of opinion leaders \cite{liu2024skepticism,gao2023s3,cai2024language,rossetti2024social,zhang2024large,liu2024tiny}. In such scenarios, the roles and relationships of agents play a critical role in ensuring realistic simulations. Initially, many studies relied on real-world data scraped from platforms to maintain consistency \cite{mou2024unveiling,gao2023s3}. However, as the scale of these simulations grew and data acquisition became more challenging, researchers began exploring the use of synthetic data \cite{yang2024oasis}. Furthermore, to accommodate the increasing demand for simulating larger numbers of agents, some studies have developed large-scale society simulation platforms \cite{gao2024agentscope,pan2024very}, employing parallel processing and other strategies to enhance simulation efficiency.

\paragraph{Recommendation Environments}
Another widely studied scenario is the recommendation environment, where these works use agents to simulate user responses in order to validate and improve recommendation algorithms~\cite{huang2023recommender,zhang2024prospect}. A key feature across these studies is the use of agents to emulate personalized behaviors such as item selection, preferences, and emotional responses, often integrating user memory and contextual factors~\cite{wang2023recagent,zhang2024agentcf,Zhang2024on}. Additionally, some approaches incorporate external knowledge or self-reflection mechanisms, allowing agents to adapt and learn from their interactions over time~\cite{wang2023recmind}. These studies collectively show how LLMs can bridge the gap between traditional recommender systems and more interactive, human-like behavior simulations, offering new ways to improve recommendation accuracy and better understand user dynamics.




\subsection{Evaluation}
\label{subsec:soc_eval}

For society simulations, the evaluation primarily focuses on the comparison between the simulation results and real-world data, with assessments centered on micro level, macro level and system level.

\paragraph{Micro-level Evaluation}
Individual simulation accuracy is key to society simulation. Therefore, micro-level evaluation of society simulation has received widespread attention. Initially, evaluations in non-real-world simulations draw on the Turing test, assessing agent behavior's resemblance to human behavior, often \textit{\underline{subjectively}} by humans or LLMs~\cite{park2023generative,liang2023leveraging,yu2024affordable}. For specific scenarios, metrics like partisan bias and human likeness index are proposed~\cite{chuang2024wisdom}. When simulations target real-world scenarios with available empirical data, automated metrics like emotion, attitude, behavior consistency, and user taste alignment can be designed for more \textit{\underline{objective}} evaluations by comparing simulation content with real-world data~\cite{gao2023s3,mou2024unveiling,Zhang2024on}.

\paragraph{Macro-level Evaluation}
Social interactions often lead to collective outcomes, so it is important to evaluate whether macro-level outcomes show patterns and trends that are consistent with the real world. For sociology and online platforms, attention is typically given to whether the scale of propagation, the distribution and trends of collective opinions and traits align with those of the real world. In addition to qualitative methods such as \textit{\underline{subjective}} evaluation~\cite{gao2023s3,Zhang2024on}, some studies have proposed quantitative metrics, such as fitted parameters, correlation coefficients and change of toxicity of community content to measure this differences \textit{\underline{objectively}}~\cite {tornberg2023simulating,liu2024skepticism,mou2024unveiling,yang2024oasis}. Similarly, in economic simulation, the evaluation of simulated economic systems depends on whether they can reproduce the most representative macroeconomic laws~\cite{li2024econagent}.

\paragraph{System-level Evaluation}
System-level evaluation is concerned with assessing the overall performance of a simulation system, irrespective of the specific content being simulated. With the growing number of agents in simulation, the focus of contemporary research has been on system efficiency and associated costs. Efficiency is assessed through various metrics, such as the time it takes to run a simulation, the resources that are utilized during the process, and how well the simulation can scale with an increasing number of agents~\cite{wang2023recagent,yang2024oasis,pan2024very}. These metrics are crucial for understanding how well the system can handle complexity and the demands of larger simulations. On the cost side, evaluations often center on the number of tokens consumed during the simulation or the financial expenditure incurred~\cite{yu2024affordable}. 
